% Supplemental section for reviewer-facing faithful Drifting reproduction.
% Keep this outside the 8-page main text unless space permits.
% For final AAAI source packaging, inline the generated table contents if
% separate \input files are disallowed by the submission source rules.

\section{Faithful Drifting Reproduction}
\label{sec:supp-faithful-drifting}

This section separates faithful reproduction of Drifting Models from the
DriftingMol-specific molecular extensions used in the main paper. The goal is
to verify that the original drift-field recipe is implemented correctly before
evaluating decoder-coupled molecular adaptations.

\paragraph{Algorithm audit.}
We audit the implementation against the original Drifting Models drift-field
construction. The checked implementation uses generated samples as negatives,
positive samples from the target QED bin, an L2 kernel, joint positive/negative
normalization, bidirectional softmax normalization, cross-multiplication of
positive and negative masses, and the anti-symmetric field $V = V^+ - V^-$.
The tensor-level audit compares the production implementation with a direct
pseudocode transcription on synthetic tensors; the current audit artifact is
\path{results/drifting_faithfulness_status.json}.

\paragraph{Strict molecular analogue.}
The strict reproduction experiments use a frozen latent-MAE feature extractor
as $\phi$, QED quantile bins as class labels, generated samples as negatives,
$N_c=64$, $N_{\mathrm{pos}}=64$, $N_{\mathrm{neg}}=64$, temperatures
$\{0.02,0.05,0.2\}$, Appendix-A.6 feature-distance/drift normalization, and
no z-diversity or decoder-coupled drift terms. Property-aware, random-feature,
and z-space controls test whether the learned feature metric is necessary.

\begin{table}[!h]
\centering
\caption{Strict faithful-Drifting core runs. These rows use the Algorithm-2
drift field without decoder-coupled drift or diversity regularization.}
\label{tab:supp-faithful-core}
\resizebox{\columnwidth}{!}{%
% Generated by scripts/collect_faithful_drifting_results.py: strict core faithfulness runs
\begin{tabular}{l c c c c c c}
\toprule
Run & Status & $\alpha$ & $\rho$ & U (\%) & MAE & Slope \\
\midrule
Strict plain-$\phi$ & complete\_pass & alpha=3.0 & 0.055 & 81.4 & 0.247 & 0.064 \\
Property-aware $\phi$ & complete\_pass & alpha=5.0 & 0.105 & 87.2 & 0.254 & 0.114 \\
Random $\phi$ & complete\_pass & alpha=5.0 & 0.122 & 97.6 & 0.222 & 0.168 \\
z-space drift & complete\_pass & alpha=5.0 & 0.127 & 97.8 & 0.221 & 0.169 \\
\bottomrule
\end{tabular}

}
\end{table}

\paragraph{Positive and negative allocation.}
The allocation sweep mirrors the original Drifting Models sample-estimation
question: whether increasing the number of positives or generated negatives
improves the drift estimate under a fixed effective batch budget. All six
allocation rows are complete. They should be read as boundary-condition
evidence: increasing $N_{\mathrm{pos}}$ from 1 through 64 improves uniqueness,
but the generated QED distribution remains concentrated near the decoder's
central QED region, so target tracking is still weak. The two negative-sample
rows produce the strongest allocation results ($\rho\approx0.076$) and recover
about 80\% uniqueness, but they remain below the pre-specified target-tracking
threshold.

\begin{table}[!h]
\centering
\caption{Positive/negative allocation sweep for faithful Drifting. Allocation
improves diversity more reliably than QED target tracking.}
\label{tab:supp-faithful-allocation}
\resizebox{\columnwidth}{!}{%
% Generated by scripts/collect_faithful_drifting_results.py: positive/negative allocation runs
\begin{tabular}{l c c c c c c}
\toprule
Run & Status & $\alpha$ & $\rho$ & U (\%) & MAE & Slope \\
\midrule
$N_{neg}=16$ & complete\_pass & alpha=5.0 & 0.076 & 79.4 & 0.246 & 0.092 \\
$N_{neg}=32$ & complete\_pass & alpha=5.0 & 0.076 & 79.9 & 0.245 & 0.092 \\
$N_{pos}=1$ & complete\_pass & alpha=3.0 & 0.028 & 68.7 & 0.251 & 0.035 \\
$N_{pos}=16$ & complete\_pass & alpha=5.0 & 0.053 & 76.7 & 0.246 & 0.065 \\
$N_{pos}=32$ & complete\_pass & alpha=5.0 & 0.062 & 79.4 & 0.244 & 0.078 \\
$N_{pos}=64$ & complete\_pass & alpha=3.0 & 0.055 & 81.4 & 0.247 & 0.064 \\
\bottomrule
\end{tabular}

}
\end{table}

\paragraph{Claim policy.}
The strict core evidence should be described conservatively: direct faithful
Drifting is reproducible in the molecular latent setting, but the observed
control is modest and remains weaker than the decoder-coupled main method.
Current allocation rows strengthen this conservative interpretation rather
than the main performance claim. The faithfulness audit is now closed for the
strict Drifting reproduction package, but the observed molecular control signal
supports a representation-limited interpretation rather than a feature-space
superiority claim.
